\section*{Block-by-block reading}

The critic $D_\phi$ is the topological dual of the generator: it consumes
a $1 \times 64 \times 64$ image (either a real training tensor or a fake
$G_\theta(z)$) and emits a single real-valued score whose difference of
expectations over $\mathbb{P}_r$ and $\mathbb{P}_g$ approximates the
Wasserstein-1 distance $W_1(\mathbb{P}_r, \mathbb{P}_g)$. The diagram
is again read top-to-bottom and uses the same colour convention as the
generator companion: \emph{green} blocks denote tensors at the network
boundary ($\mathbf{x}$ and $D_\phi(\mathbf{x})$), \emph{blue} blocks
denote internal feature maps $\mathbf{h}^{(\ell)}$ that the critic
materializes and discards during a forward pass, and \emph{amber} blocks
denote the parameterised operators. The side labels record the spatial
resolution at each layer; in the critic, that resolution is \emph{halved}
at every stage, mirroring the generator's doubling pattern. Below we walk
through every block in the order in which the forward pass visits them.
